\documentclass[a4paper,oneside,12pt]{extreport}

\usepackage[utf8]{inputenc}
\usepackage[english]{babel}
\usepackage{amsmath}
\usepackage{amssymb}
\usepackage{amsthm}
\usepackage{enumitem}
\usepackage{geometry}
\usepackage{mathtools}


\geometry
{
lmargin=15mm,
rmargin=15mm,
tmargin=0mm,
bmargin=5mm,
includeheadfoot
}
\pagestyle{empty}
\frenchspacing


\theoremstyle{plain}
\newtheorem{theorem}{Theorem}
\newtheorem{lemma}[theorem]{Lemma}
\newtheorem{proposition}[theorem]{Proposition}
\newtheorem{corollary}[theorem]{Corollary}

\theoremstyle{definition}
\newtheorem{definition}[theorem]{Definition}
\newtheorem{example}[theorem]{Example}

\theoremstyle{remark}
\newtheorem{remark}[theorem]{Remark}


\newcommand*{\ve}{\varepsilon}
\newcommand*{\mbP}{\mathbb{P}}
\newcommand*{\mbQ}{\mathbb{Q}}
\newcommand*{\mbR}{\mathbb{R}}
\newcommand*{\mcF}{\mathcal{F}}
\newcommand*{\mcN}{\mathcal{N}}
\newcommand*{\mcP}{\mathcal{P}}

\newcommand*{\ges}{\geqslant}
\newcommand*{\les}{\leqslant}
\newcommand*{\indf}{1\negmedspace\mathrm{I}}
\newcommand*{\E}{\mathbb{E}}
\newcommand*{\LV}{\text{LV}}
\newcommand*{\SV}{\text{SV}}

\DeclarePairedDelimiter{\abs}{\lvert}{\rvert}






\begin{document}
\begin{center}
{\large\bfseries Notes on $p$-variation}\\[6pt]
{\small Technical notes, 2026}
\end{center}
\bigskip


For an arbitrary function $f\colon [0,T]\to\mbR$ and $p>0$ define
\begin{gather*}
\LV_p(f)=\text{LV}_p(f,[0,T])=\lim_{\Delta(\mcP)\to 0+} S_p(f;\mcP),\\
\SV_p(f)=\text{SV}_p(f,[0,T])=\sup_{\mcP} S_p(f;\mcP),
\end{gather*}
where
\[
S_p(f;\mcP)=\sum_{i=1}^n (f(t_i)-f(t_{i-1}))^p.
\]
Note that the limit may not exist, whereas the supremum always exists, even though it can infinite.


\begin{proposition}
It holds
\[
\LV_p(f)\les\SV_p(f).
\]
\end{proposition}


\begin{proof}
This follows immediately from the fact that for any partition $\mcP$ it holds
\[
S_p(f;\mcP)\les\sup_{\mcP'} S_p(f;\mcP')=\SV_p(f).
\]
\end{proof}


\begin{example}
Consider the function
\[
f(t)=
\begin{cases}
0,&\text{if $0\les t<1/2$,}\\
h,&\text{if $1/2\les t\les 1$.}
\end{cases}
\]
where $h$ is an arbitrary real number. It is clear that
\[
\LV_p(f)=\SV_p(f)=\abs{h}^p.
\]
Depending on the value of $\abs{h}$, we have the following situations:
\begin{itemize}
\item
if $\abs{h}<1$, then both $\LV_p(f)$ and $\SV_p(f)$ are decreasing in $p$,

\item
if $\abs{h}=1$, then both $\LV_p(f)$ and $\SV_p(f)$ are constant in $p$,

\item
if $\abs{h}>1$, then both $\LV_p(f)$ and $\SV_p(f)$ are increasing in $p$.
\end{itemize}
\end{example}


\begin{example}
Consider the function
\[
f(t)=
\begin{cases}
0,&\text{if $0\les t<1/3$,}\\
h_1,&\text{if $1/3\les t<2/3$,}\\
h_1+h_2,&\text{if $2/3\les t\les 1$,}
\end{cases}
\]
where $h_1,h_2\ges 0$ are arbitrary real numbers. It is clear that
\[
\LV_p(f)=h_1^p+h_2^p.
\]

Now take an arbitrary partition $\mcP$ of the interval $[0,T]$ and consider two cases:
\begin{itemize}
\item
if the interval $[1/3,2/3)$ does not contain any points of the partition, then
\[
S_p(f;\mcP)=(h_1+h_2)^p.
\]

\item
if the interval $[1/3,2/3)$ contains at least one point of the partition, then
\[
S_p(f;\mcP)=h_1^p+h_2^p.
\]
\end{itemize}
Therefore,
\[
\SV_p(f)=\max\{(h_1+h_2)^p,h_1^p+h_2^p\}.
\]
Note that, generally speaking, we do not know which of the two values is larger. For instance, if $h_1=h_2=1/2$ and $p=2$, we have
\[
(h_1+h_2)^p=1>1/2=h_1^p+h_2^p.
\]
However, if $h_1=h_2=1/2$ and $p=1/2$, we have
\[
(h_1+h_2)^p=1<\sqrt{2}=h_1^p+h_2^p.
\]


\begin{lemma}
For any $x\ges 0$ it holds
\begin{equation*}
\begin{aligned}
1+x^p&\ges (1+x)^p,\qquad\text{if $0<p<1$,}\\
1+x^p&\les (1+x)^p,\qquad\text{if $p\ges 1$.}
\end{aligned}
\end{equation*}
\end{lemma}


\begin{proof}
Consider the function
\[
g(x)=(1+x)^p-1-x^p.
\]
For $0<p<1$ we have
\[
g'(x)=p(1+x)^{p-1}-px^{p-1}<0,
\]
and for $p\ges 1$ we have
\[
g'(x)=p(1+x)^{p-1}-px^{p-1}\ges 0.
\]

Therefore, for any $x\ges 0$, we get
\begin{gather*}
\begin{aligned}
g(x)&\les g(0)=0,\qquad\text{if $0<p<1$,}\\
g(x)&\ges g(0)=0,\qquad\text{if $p\ges 1$,}
\end{aligned}
\end{gather*}
which implies the required inequalities.
\end{proof}


Using this lemma, we get that if $p\ges 1$, then
\[
(h_1+h_2)^p=h_1^p\left(1+\frac{h_2}{h_1}\right)^p\ges h_1^p\left(1+\frac{h_2^p}{h_1^p}\right)=h_1^p+h_2^p,
\]
and if $0<p<1$, then
\[
(h_1+h_2)^p=h_1^p\left(1+\frac{h_2}{h_1}\right)^p\les h_1^p\left(1+\frac{h_2^p}{h_1^p}\right)=h_1^p+h_2^p.
\]
Therefore,
\[
\SV_p(f)=
\begin{cases}
h_1^p+h_2^p,&\text{if $0<p<1$,}\\
(h_1+h_2)^p,&\text{if $p\ges 1$.}
\end{cases}
\]
\end{example}
\end{document}